% ============================================================
%  chapters/ch11-trajectories-admissibility.tex
% ============================================================

\chapter{Trajectories and Admissibility}

\chapterepigraph{%
  A gesture that cannot be admitted\\
  is not a gesture — it is noise.\\
  Admissibility is what turns motion into meaning.
}{Flyxion, \textit{Semantic Infrastructure}}

\sidenote{App.\,F for gesture\\manifolds; App.\,H\\for admissibility\\functionals}

Chapter~10 established that gestures are trajectories in constrained
manifolds and that symbols are projections of gesture equivalence classes.
This chapter develops the mathematical structure of the gesture manifold
more carefully, focusing on the \emph{admissibility conditions} that
separate meaningful gesture from random motion.

\section{Gesture Manifolds}

The gesture space $\gesturesp = \mathrm{Path}([0,T], X)$ of all possible
hand trajectories is enormous — a function space of infinite dimension.
Most of this space is physically inaccessible, kinematically impossible,
or neurologically non-executable. The gesture manifold $M_{gestureidx} \subset
\gesturesp$ is the small submanifold of trajectories that real hands can
actually produce.

\begin{definition}{Biomechanical Constraint Manifold}{biomech-mfld}
  The \emph{biomechanical constraint manifold} $M_\text{bio} \subset
  \gesturesp$ is the set of trajectories satisfying:
  \begin{enumerate}
    \item \textbf{Range:} each joint angle $\theta_i$ satisfies
          $\theta_i^{\min} \leq \theta_i \leq \theta_i^{\max}$.
    \item \textbf{Velocity:} angular velocity $|\dot{\theta}_i| \leq v_i^{\max}$.
    \item \textbf{Coordination:} adjacent finger joints move in
          coordinated patterns (flexion synergies, extension synergies).
    \item \textbf{Continuity:} trajectories are smooth ($C^2$ at minimum).
  \end{enumerate}
\end{definition}

Within $M_\text{bio}$ lies a further restriction: the \emph{skill
manifold} $M_\text{skill} \subset M_\text{bio}$ of trajectories that
a particular person, with their particular training and motor history,
can execute reliably. A concert pianist's skill manifold includes
trajectories that a novice's does not; a surgeon's includes trajectories
that are biomechanically possible but not executable without years of
practice.

\section{Constraint Spaces on Trajectory Space}

The admissibility framework of Chapter~1 and Appendix~H applies to
trajectory spaces directly.

\begin{definition}{Gesture Constraint Space}{gesture-constraint}
  A \emph{gesture constraint space} is a triple
  $\mathcal{C}_{gestureidx} = (\gesturesp, A_{gestureidx}, \kappa_{gestureidx})$
  where the admissible region $A_{gestureidx} = \kappa_{gestureidx}^{-1}(0)$
  is the set of gestures that satisfy all biomechanical, skill, and
  contextual constraints.
\end{definition}

The constraint functional $\kappa_{gestureidx} : \gesturesp \to \mathbb{R}_{\geq 0}$
measures total constraint violation. A typical decomposition:
\[
  \kappa_{gestureidx}(\gamma) = \kappa_\text{bio}(\gamma) + \kappa_\text{skill}(\gamma)
  + \kappa_\text{context}(\gamma)
\]
where each term measures violation of biomechanical, skill, and contextual
constraints respectively.

\section{Motor Equivalence Classes Revisited}

Two admissible gestures are motor-equivalent when they differ only within
the fiber of the symbolic projection — when their difference is invisible
at the symbolic output level.

\begin{definition}{Admissible Motor Equivalence}{adm-motor-equiv}
  Two gestures $\gamma_1, \gamma_2 \in A_{gestureidx}$ are
  \emph{admissibly motor-equivalent}, written $\gamma_1 \approx_m \gamma_2$,
  if $\Phi(\hat{C}(\gamma_1)) = \Phi(\hat{C}(\gamma_2))$ — they produce
  the same symbolic output after canonicalization.
\end{definition}

The admissibility condition ensures we only identify gestures that are
genuinely executable, not merely gestures that would produce the same
symbol if they could be executed.

\begin{example}
  Writing the letter ``e'' with the dominant hand and the non-dominant
  hand both produce the symbol ``e'', so $\gamma_R \approx_m \gamma_L$.
  But $\gamma_L$ is often in the skill manifold for right-handed writers
  only under unusual conditions — it is admissible but not in $M_\text{skill}$.
  Admissibility conditions prevent us from identifying these as equivalent
  in normal writing contexts.
\end{example}

\section{Canonical Gestures}

The canonical gesture is the minimum-energy, maximum-reliability
representative of an equivalence class.

\begin{definition}{Canonical Gesture}{canonical-gesture}
  The \emph{canonical gesture} $\hat{\gamma}$ for equivalence class
  $[\gamma]_{\approx_m}$ is the solution to:
  \[
    \hat{\gamma} = \arg\min_{\gamma' \approx_m \gamma} E(\gamma')
    \quad \text{subject to} \quad \gamma' \in A_{gestureidx}
  \]
  where $E(\gamma') = \int_0^T \|\dot{\gamma}'(t)\|^2\,dt$ is the
  kinetic energy.
\end{definition}

\begin{theorem}{Canonical Gesture Existence}{canonical-existence}
  Under mild regularity conditions on $A_{gestureidx}$
  (weak compactness of the admissible set and lower semicontinuity of $E$),
  every non-empty equivalence class has a canonical gesture.
\end{theorem}

\begin{proof}[Sketch]
  By the direct method of the calculus of variations: the energy
  functional $E$ is coercive and lower semicontinuous on the weakly
  compact admissible set; therefore it attains its infimum.
\end{proof}

\section{Accessibility of Gesture Continuations}

The accessibility entropy of a gesture $\gamma$ at time $t$ measures
how many admissible continuations are available:
\[
  S(\gamma, t) = \log \mathrm{Vol}\bigl(\{\gamma' \in A_{gestureidx} :
  \gamma'|_{[0,t]} = \gamma|_{[0,t]}\}\bigr).
\]
This measures the remaining freedom given the past trajectory — how
many ways the gesture can continue.

\begin{exercise}
  Show that the accessibility entropy $S(\gamma, t)$ is non-increasing
  in $t$ along any fixed trajectory $\gamma$. What does equality
  $S(\gamma, t+\delta) = S(\gamma, t)$ imply about the gesture at time $t$?
\end{exercise}

\begin{exercise}
  Define the \emph{gesture manifold curvature} at a point
  $\gamma \in M_{gestureidx}$ as the second fundamental form of the
  embedding $M_{gestureidx} \hookrightarrow \gesturesp$. Interpret
  high curvature geometrically. What kind of gesture motions
  correspond to regions of high curvature?
\end{exercise}

\begin{exercise}[title={Open problem}]
  The canonical gesture minimizes kinetic energy among equivalent
  gestures. Is this the same as maximizing reliability (probability
  of correct recognition under noise)? If not, construct an example
  where the minimum-energy gesture is less reliable than a higher-energy
  alternative. What modified optimality criterion resolves the tension?
\end{exercise}
